\noindent\textit{2008 manuscript.}

\section{Phase Transformations}

\subsection*{Vapor Pressure Equation}

\subsubsection*{Derivation of the Coexistence Curve $p$ versus $\tau$}

Suppose $\mu_g(p_0,\tau_0) = \mu_l(p_0,\tau_0)$ and $\mu_g(p_0 + dp, \tau_0 + d\tau) = \mu_l(p_0 + dp, \tau_0 + d\tau)$
\[
\Longrightarrow \mu_g(p_0,\tau_0) + \left( \frac{ \partial \mu_g }{ \partial p} \right)_{\tau} dp + \left( \frac{ \partial \mu_g }{ \partial \tau} \right)_{p} d\tau + \dots = \mu_l(p_0,\tau_0) + \left( \frac{ \partial \mu_l }{ \partial p} \right)_{\tau} dp + \left( \frac{ \partial \mu_l }{ \partial \tau} \right)_{p} d\tau + \dots
\]

\[
\frac{dp}{d\tau} = \frac{ \left( \frac{ \partial \mu_l }{ \partial \tau} \right)_p - \left( \frac{ \partial \mu_g }{ \partial \tau} \right)_p }{ \left( \frac{ \partial \mu_g }{ \partial \tau} \right)_{\tau} - \left( \frac{ \partial \mu_l }{ \partial p} \right)_{\tau} }
\]




$\begin{aligned}
  & v \equiv V/N \\ 
  & s \equiv \sigma /N
\end{aligned}$

\[
\begin{aligned}
  & \frac{1}{N} \left( \frac{ \partial G}{ \partial p } \right)_{N,\tau} = v = \left( \frac{ \partial \mu }{ \partial p } \right)_{\tau} \\ 
  & \frac{1}{N} \left( \frac{ \partial G}{ \partial \tau} \right)_{N,p} = \frac{-\sigma}{N} = -s = \left( \frac{ \partial \mu }{ \partial \tau} \right)_p
\end{aligned}
\]


$s_g-s_l$ related directly to quantity of heat that must be added to system to transfer 1 molecule reversibly from liquid to gas, while temperature constant.  (if heat isn't added to system from outside in the process, temperature will decrease when molecule transferred to gas)
\begin{equation}
  L \equiv \tau(s_g - s_l)
\end{equation}
\textbf{latent heat of vaporization}

\begin{equation}
  \Delta v = v_g - v_l
\end{equation}
change of volume when 1 molecule transferred from liquid to gas

\begin{equation}
  \frac{dp}{d \tau} = \frac{ L }{ \tau \Delta v}
\end{equation}
is the \textbf{ Clausius-Clapeyron equation} or \textbf{vapor pressure equation}.  

Assumptions:
\begin{enumerate}
  \item[(a)] if assume $v_g \gg v_l$, then $\Delta v \cong v_g = V_g/N_g$ at $p=1$ atm (at atmospheric pressure), $v_g/v_l \approx 10^3$
  \item[(b)] ideal gas law $pV_g = N_g \tau \Longrightarrow \Delta v \cong \tau/p$
\end{enumerate}
\[
\frac{dp}{d\tau} = \frac{L}{\tau^2}p \text{ or } \frac{d}{d\tau} \ln{p} = \frac{L}{\tau^2}
\]
$L \equiv $ latent heat per molecule

If latent heat $L$ independent of temperature over temperature range of interest, $L = L_0$ \\
\phantom{\quad \, } $\int \frac{dp}{p} = L_0 \int \frac{d\tau}{\tau^2}$ whence $\ln{p} = -L_0/\tau + \text{const}$ or $p(\tau) = p_0\exp{ (-L_0/\tau)}$



Consider $\frac{L_0}{\tau}$, $L_0$, latent heat of vaporization of \emph{1 molecule}.  

Now compare it with $\frac{L_0}{RT}$, \\
where $R \equiv N_0 k_B \equiv $ gas constant.  

$N_0 = 6.02205 \times 10^{23} \frac{1}{ \text{mol}}$ \\
$k_B = 1.38066 \times 10^{-23} \, J/K$ (usual Boltzmann constant) \\
$R = 8.31441 \, J/K\cdot \text{mol}$ 

So for 
\[
\frac{L_0}{RT} = \frac{L_0}{N_0 k_B T}
\]
Then that implies
\[
\frac{1}{ 1 \text{ particle} } \frac{1}{N_0} = \frac{1}{ 1 \text{ particle } \frac{1}{ \text{ mol }} }
\]
which implies that $L_0$ now is the latent heat of vaporization of a \emph{mole (mol)} of the gas.  



\subsubsection*{Latent heat and enthalpy}

crossing the coexistence curve,
\[
\tau d\sigma = dU + pdV - (\mu_g - \mu_l) dN
\]
if liquid $\to$ gas, $\mu_l > \mu_g$ (or $0 > \mu_g - \mu_l$)

on coexistence curve, $\mu_g = \mu_l$
\[
L = \tau d\sigma(\dot{\gamma}) = dU(\dot{\gamma}) + pdV(\dot{\gamma}) = (dU+ pdV)(\dot{\gamma}) + 0 = (dU + pdV)(\dot{\gamma}) + Vdp(\dot{\gamma}) = dH(\dot{\gamma}) = H_g -H_l
\]
where $dp(\dot{\gamma}) = 0$ since $\gamma$ is a path s.t. pressure $p$ is constant.  

Recall that $\begin{gathered} \quad \\
  H = U + pV \\
  dH = dU+ pdV + Vdp \end{gathered}$.  


$C_p = \tau \left( \frac{ \partial \sigma }{ \partial \tau} \right)_p = \tau\sigma(\tau,p)(\dot{\gamma}) = (dU(\tau,p)(\dot{\gamma}) + pdV(\tau,p)(\dot{\gamma}) ) = \left( \frac{ \partial H}{ \partial \tau} \right)_p$

Then $H = \int C_p d\tau = \int_{\gamma} C_p d\tau$

\textbf{Example: Model system for gas-solid equilibrium} (pp.285 Ch. 10 ``Phase transformations'', Kittel and Kroemer (1980) \cite{CKittelHKroemer1980})

imagine solid of $N$ atoms \\
each atom bound as harmonic oscillator of frequency $\omega$ \\
binding energy of each atom in ground state $\epsilon_0$, i.e. energy of atom in ground state is $-\epsilon_0$ \\
energy of a single oscillator $\epsilon = n\hbar \omega - \epsilon_0$, $n\in \mathbb{Z}^+$ 

$Z_s \equiv $ partition function of a single oscillator in solid, $\beta \equiv 1/\tau$
\[
Z_s = \sum_n \exp{ [ -\beta (n\hbar \omega - \epsilon_0) ]} = \epsilon^{\beta \epsilon_0} \sum_n (e^{-\beta \omega} )^n = \frac{ e^{\beta \epsilon_0}}{ 1 - e^{-\beta \omega }}
\]

Free energy $F_s = U_s - \tau \sigma_s = -\tau \ln{Z_s}$

Gibbs free energy in the solid, per atom 
\[
G_s = U_s - \tau \sigma_s + pv_s = F_s + pV_s = \mu_s
\]


Supposing $v_s \ll v_g$ (volume $v_s$ per atom in solid phase much smaller than volume $v_g$ per atom in gas phase)
\[
\begin{gathered}
  \mu_s = F_s + pv_s \approx F_s \\ 
  \lambda_s := \exp{ (\beta \mu_s)} \simeq \exp{ (\beta F_s)} = \exp{ (-\ln{ Z_s }) } = \frac{1}{Z_s} = \exp{ (-\beta \epsilon_0 )} (1- \exp{ (-\beta \omega )} )
\end{gathered}
\]
use ideal gas approximation to describe gas phase, and spin of atom to be $0$.  From Ch. 6,
\[
\lambda_g = \frac{n}{n_Q} = \frac{p}{\tau n_Q} = \frac{p}{\tau} \left( \frac{2\pi \hbar^2 }{ M \tau} \right)^{3/2}
\]

gas in equilibrium with solid with $\lambda_g = \lambda_s$, or 
\[
\begin{gathered}
  \frac{p}{ \tau n_Q} = \lambda_g = \lambda_s = e^{-\beta \epsilon_0} ( 1-  e^{-\beta \omega } ) \Longrightarrow p = \left( \frac{M}{ 2\pi \hbar^2 } \right)^{3/2} \tau^{5/2} \exp{ (-\epsilon_0/\tau) } (1- \exp{ ( -\hbar \omega / \tau ) } )
\end{gathered}
\]


\subsection*{Van Der Waals Equation of State}


\[
F(\text{ideal gas}) = -N\tau [ \ln{ (n_Q/n) } + 1 ]
\]
$n:= \frac{N}{V} = \frac{N}{V-Nb}$ (hardcore repulsion at short distances)

\subsubsection*{Mean Field Method}

Let $\varphi(r) := $ potential energy of interaction of 2 atoms separated by distance $r$.

average value of total interaction of all other atoms on atom at $r=0$
\[
\int_0^{\infty} dV \varphi(r)n = n \int_b^{\infty} dV \varphi(r) = -2na \quad \, (2 \text{ factor is convention})
\]

excluded hardcore sphere of volume $b$ from volume of integration. \\
assumed mean value of $n$, concentration $n$ constant throughout volume accessible to gas molecules. \\
\phantom{ \quad \, } ignored increase of concentration in regions of strong attractive potential energy i.e. \\
\phantom{ \quad \quad \, } mean field method meglects correlations between interacting molecules.

$\Delta F \simeq \Delta U = \frac{-1}{2} (2 N na) = -N^2 a/ V $ \\
\phantom{\quad \, } where $\frac{1}{2} N(N-1) \simeq \frac{1}{2} N^2 $ ($N$ large)

\[
\Longrightarrow F(\text{vdW}) = -N\tau ( \ln{ [n_Q (V-Nb)/N ]} + 1) - N^2 a /V
\]
\[
p = -\left( \frac{ \partial F}{ \partial V} \right)_{\tau,N} = \frac{N\tau}{ V-Nb} - \frac{N^2 a}{V^2} \text{ or } \left( p + \frac{N^2 a}{V^2} \right)(V-Nb) = N\tau
\]

\subsubsection*{Critical Points for the Van der Waals Gas}

\[
\begin{aligned}
  & p_c = \frac{a}{27 b^2} \\ 
  & V_c = 3Nb \\ 
  & \tau_c = \frac{8a}{ 27b}
\end{aligned}
\]

\[
\begin{aligned}
  & \widehat{p} := \frac{p}{p_c} \\ 
  & \widehat{V} := \frac{V}{V_c} \\ 
  & \widehat{\tau} := \frac{ \tau}{ \tau_c}
\end{aligned}
\]
since
\[
\begin{gathered}
\begin{aligned}
  & \left( \frac{p}{p_c} + \frac{3}{ (V/V_c)^2 } \right)\left( \frac{V}{V_c} - \frac{1}{3} \right) = \frac{8 \tau }{ 3\tau_c} \\
  & \left( \frac{p}{p_c} + \frac{3}{ (V/V_c)^2} \right) = \frac{1}{p_c} ( p + \frac{3a/27b^2}{ V^2/9N^2 b^2 } ) = \frac{1}{p_c} ( p + \frac{N^2 a}{V^2} ) 
\end{aligned} \quad \quad \, 
\begin{aligned}
  & \frac{8\tau}{3\tau_c} = \frac{8 \tau}{ 3 \frac{ 8a}{27b} } = \frac{9\tau}{a/b} \\ 
  & \frac{1}{p_c V_c} = \frac{1}{ \frac{a}{27 b^2} 3Nb } = \frac{9}{a/b}
\end{aligned} \\
\left( \frac{V}{V_c} - \frac{1}{3} \right) = \frac{1}{V_0}(V-Nb)
\end{gathered}
\]

\[
\Longrightarrow \left( \widehat{p} + \frac{3}{ \widehat{V}^2 } \right)\left( \widehat{V} - \frac{1}{3} \right) = \frac{8}{3} \widehat{\tau} \text{ or } \widehat{p} = \frac{ \frac{8}{3} \widehat{\tau} }{ \widehat{V} - \frac{1}{3} } - \frac{3}{ \widehat{V}^2}
\]

at $\widehat{p} =1$, $\widehat{V}=1$, $\widehat{\tau}=1$, $\widehat{\tau}$ constant, $\left( \frac{ \partial \widehat{p}}{ \partial \widehat{V}} \right)_{\widehat{\tau}} =0$, $\left( \frac{ \partial^2 \widehat{p}}{ \partial \widehat{V}^2} \right)_{\widehat{\tau}} =0$ \\
Above $\tau_c$, no phase separation exists.  

\subsubsection*{Gibbs Free Energy of the Van der Waals Gas} $G= F+pV$
\[
G(\tau,V,N) = \frac{ N \tau V}{ V - Nb} - \frac{2N^2 a}{V} - N\tau ( \ln{ [n_Q(V-Nb)/N] } + 1)
\]

Recall $dG = \mu dN - \sigma d\tau + Vdp$
\[
\Longrightarrow \left( \frac{ \partial G}{ \partial N} \right)_{\tau,p} = \mu = \mu(p,\tau) \text{ or } G(N,p,\tau) = N\mu(p,\tau)
\]
cf. (9.13) of Kittel and Kroemer (1980) \cite{CKittelHKroemer1980}

$G(\tau,p,N)$ solved numerically.

on p-V diagram, \\
\phantom{ \, } coexistence line: 
\phantom{ \, \quad } consider $dG = -\sigma d\tau + Vdp + \mu dN$ \\
\phantom{ \, \quad \, } at constant $\tau$, constant total number of particles, $N$, \\
$dG(\dot{\gamma}) = Vdp(\dot{\gamma}) = G_g - G_l = \int V dp$


if $\int V dp =0$, $G_p =G_l$, $G_g(\tau,p) = G_l(\tau,p)$, so $\mu_g(\tau,p) = \mu_l(\tau,p)$ along horizontal coexistence line

\subsubsection*{Nucleation}

Let $\Delta \mu = \mu_g - \mu_l$ chemical potential difference between \\
vapor surrounding small liquid droplet and \\
liquid in bulk (infinitely large drop) 

surface free energy of liquid drop is positive and tends to increase free energy of liquid \\
for small drop radii, surface can be dominant and drop can be unstable with respect to gas 

$n_l \equiv $ concentration of molecules in liquid 
\[
\Delta G = G_l -G_g = - \left( \frac{4\pi}{3} \right) R^3 n_l \Delta \mu + 4\pi R^2 \gamma 
\]
where $\gamma$ surface free energy per limit area.  

\[
\frac{d \Delta G}{dR} = 0 = -4\pi R^2 n_l \Delta \mu + 8 \pi R \gamma \text{ or } R_c = \frac{2\gamma }{ n_l \Delta \mu}
\]
$R_c \equiv $ critical radius for nucleation \\
when $R< R_c$, $\Delta G >0$, so drop will evaporate (to gas) spontaneously \\
when $R >R_c$, $\Delta G <0$, drop will tend to grow spontaneously

\[
(\Delta G)_c = \frac{16 \pi}{3} [ \frac{ \gamma^3}{ n_l^2 (\Delta \mu)^2 } ]
\]
where $(\Delta G)_c = \Delta G(R= R_c)$

assume vapor behaves like ideal gas: 
\[
\Delta \mu = \tau \ln{ \left( \frac{ p}{ p_{\text{eq}} } \right) }
\]
Eq. (12a), chemical potential of the ideal gas, $\mu = \tau \ln{ \left( \frac{n}{n_Q} \right)} = \tau \ln{ \left( \frac{p}{ p_{\text{eq}}} \right)}$ \\
\phantom{\quad \, } $p \equiv $ vapor pressure in gas phase \\ 
\phantom{\quad \, } $p_{\text{eq}} \equiv $ equilibrium vapor pressure of bulk liquid $(R \to \infty)$ \\
use $\gamma = 72 \, \frac{\text{erg}}{\text{cm}^2}$, $p =1.1 p_{\text{eq}}$, $\tau = 300 \, K$ ($H_2O$), $R_c = \frac{2\gamma}{n_l \Delta \mu } = 1\times 10^{-6}$ cm.  



\subsection*{Problems}

\problemhead{1} \textbf{Entropy, energy, and enthalpy of van der Waals gas}.  
\begin{enumerate}
\item[(a)]
Recall the generalized thermodynamic identity; starting with $\sigma = \sigma(U,V,N) \in C^{\infty}(\Sigma) = C^{\infty}((U,V,N))$, the identity is $dU = \tau d\sigma - pdV + \mu dN$.  

By definition, $F=U-\tau \sigma$, so that 
\[
dF = dU- d\tau \sigma -\tau d\sigma = - pdV + \mu dN - \sigma d\tau 
\]
Then clearly, $\left( \frac{ \partial F}{ \partial \tau} \right)_{V,N} = -\sigma$.  



Now 
\[
F(\text{vdW}) = -N\tau ( \ln{ [n_Q (V-Nb)/N ]} + 1) - N^2 a /V
\]
as was postulated about the hardcore repulsion inside a molecule, making $V$ into $V-Nb$.  

Remembering that quantum concentration does depend on $\tau$: $n_Q \equiv \left( \frac{ M \tau}{ 2\pi \hbar^2 } \right)^{3/2}$


Consider 
\[
\left( \ln{ ( \left( \frac{M \tau}{ 2\pi \hbar^2} \right)^{3/2} \frac{ (V-Nb)}{ N} )} + 1 \right) = \left( \frac{3}{2} \ln{\tau} + \ln{ \left( \left( \frac{M}{2\pi \hbar^2} \right)^{3/2} \frac{(V-Nb)}{N} \right) } + 1 \right) \xrightarrow{ \frac{ \partial }{\partial \tau} } \frac{3}{2\tau}
\]
Thus
\[
\sigma = - \left( \frac{ \partial F}{ \partial \tau } \right)_{V,N} = N \left( \ln{ \left( \frac{n_Q(V-Nb)}{N} \right) } +1 \right) + N \frac{3}{2} = N \left( \ln{ \left( \frac{n_Q(V-Nb)}{N} \right) } + \frac{5}{2} \right)
\]

\item[(b)] Recall that $F:= U-\tau \sigma$ or $U = F+\tau \sigma = F- \tau \left( \frac{ \partial F}{ \partial \tau} \right)_{V,N} = -\tau^2 \left( \frac{ \partial }{ \partial \tau} \left( \frac{F}{\tau} \right) \right)_{V,N}$  
\[
 U = -\tau^2 \frac{3}{2\tau}(-N) + \frac{N^2a }{V} = \boxed{ \frac{3N\tau}{2} + \frac{N^2a }{V} = U }
\]
\item[(c)] Give to first order in the van der Waals correction terms $a,b$.  

\[
p = \frac{N\tau}{V-Nb} - \frac{N^2a}{V^2} = \frac{N\tau}{ V(1- Nb/V)} - \frac{N^2 a}{V^2} = \frac{N\tau}{V} ( 1 + \frac{Nb}{V} ) - \frac{N^2 a}{V^2}
\]
where $Nb \ll V$.  

Thus,
\[
\boxed{ H(\tau,V) := U + pV = \frac{5}{2} N\tau  + \frac{N^2 b\tau}{V} - \frac{2N^2 a}{V} }
\]
\end{enumerate}


\problemhead{2} \textbf{Calculation of $dT/dp$ for water}

Using the \textbf{Clausius-Clapeyron equation} or \textbf{vapor pressure equation}, 
\[
\frac{dp}{d\tau} = \frac{L}{\tau \Delta v}
\]

assume $\Delta v := v_g - v_l \approx v_g \equiv \frac{V_g}{N_g}$ (i.e. the volume of the gas is much bigger than the volume of the liquid), and $pV_g = N_g \tau$ or $\frac{V_g}{N_g} = \frac{\tau}{p}$ (i.e. water vapor acts like an ideal gas)

\[
\frac{L}{ \tau \Delta v} \approx \frac{L}{ \tau^2 p}
\] 

Looking up fundamental constants and unit conversions \emph{manually} appears, to me at least, in the 21st century, to be an obsolete exercise (``an exercise in insanity?'').  I would like to propose presenting the fundamental constants and unit conversions, sourced from the NIST (National Institute of Standards and Technology) website, and using Python library \verb|pandas|, as panda dataframe objects, to make the data accessible and computable, and persisting them or databasing the values as panda dataframe objects.  

Open up \verb|thermo.py|, which relies on the \verb|Physique| repository (I'll put it up on my \verb|github|, \verb|ernestyalumni|).  

\begin{lstlisting}
N_A = FundConst[ FundConst["Quantity"].str.contains("Avogadro")].loc[42,:]
k_BOLTZ = FundConst[ FundConst["Quantity"].str.contains("Boltzmann") ].loc[49,:]
P_frmATM = conv[ conv["Toconvertfrom"].str.contains("atm") ].loc[15,:] # 1 atm to Pascal conversion for pressure

...

## 2. Calculation of $dT/dp$ for water

L = 2260 # J g^{-1}
boilingwatertemp_K = KCconv.subs(T_C,100).rhs # room temperature in Kelvin
tau_boilingwater = boilingwatertemp_K * k_BOLTZ.Value
P = 1 # atm

dpdtau_1002 = L*18.0153/float(N_A.Value)/( tau_boilingwater**2 )
dtaudp_1002 = 1./ dpdtau_1002 # in J/Pascal
dTdp_1002 = dtaudp_1002 / (k_BOLTZ.Value )  # 28.4348535111262 K/atm
\end{lstlisting}

$\boxed{ \frac{dT}{dp} = 28.43 K/\text{atm} }$


\problemhead{3} \textbf{Heat of vaporization of ice}

Given
\[
\begin{aligned}
  & P(\ce{H2O} \text{ vapor} )(T = -2^{\circ}C) = 3.88 \, mm \, Hg \\ 
  & P(\ce{H2O} \text{ vapor})(T = 0^{\circ}C) = 4.58 \, mm \, Hg 
\end{aligned}
\]

Use the \textbf{Clausius-Clapeyron equation} or \textbf{vapor pressure equation} (for coexistence equilibrium).  

\[
\begin{gathered}
  \frac{dp}{d\tau} = \frac{L}{ \tau \Delta v} \approx \frac{L}{ \tau^2/p} \\ 
  \text{ with } \frac{dp}{d\tau} \approx \frac{\Delta p}{ \Delta \tau} = \frac{ p(\tau_2) - p(\tau_1) }{ \tau_2- \tau_1 }
\end{gathered}
\]

Now idealizing water vapor as a perfect ideal gas, $p = \frac{N}{V}\tau$, so $p$ linear in $\tau$, i.e. $p=p(\tau)$.  

Note that while $p$ is linear in $\tau$, the given temperature values correspond to $T$, in $K$, and so one would need this form of the slope equation $y=mx+b$, where $m=\frac{N}{V}$ times a bunch of conversion factors, and one must also include a nonzero $b$ constant, the result of choice of units (celsius versus energy units).  

Thus, $L = \frac{dp}{d\tau} \frac{\tau^2}{p}$, with $L$ being the latent heat of vaporization per particle.  

Note, I treat Avogadro's number as a units conversion, $1 = N_A = 6.02205 \times 10^{23} \, \frac{ \text{ particles }}{ \text{ mol } }$.  So in the end, for the desired $J/\text{mol}$, multiply $L$ by $N_A$, $LN_A$.  

\begin{lstlisting}
L1003 = ( 4.58 - 3.88 )/(0 - (-2.) )*((KCconv.subs(T_C,1.).rhs)**2)/( (4.58-3.88)/(0 - (-2)) * ( 1.) + 3.88 ) * 
k_BOLTZ.Value*N_A.Value  # 51705.6757640485 
\end{lstlisting}

\[
\boxed{L = 5.17 \times 10^4 \, J/\text{mol} }
\]

\problemhead{4} \textbf{Gas-solid equilibrium}

\begin{enumerate}
\item[(a)] Let's review the quantum harmonic oscillator (qho).  

Recall the Hamiltonian in $d$ dimensions
\[
\widehat{H} = \frac{ \widehat{p}^2}{2m} + \frac{1}{2} m \omega^2 \widehat{x}^2
\]
where $\widehat{p} = -i \frac{ \partial }{ \partial x^i} \mathbf{e}_i$.  

Suppose for the Hilbert space $\mathcal{H}$, that $\mathcal{H} = \mathcal{H}_{x^1} \otimes \mathcal{H}_{x^2} \otimes \mathcal{H}_{x^3}$, 
\[
| \psi \rangle \in \mathcal{H} = \mathcal{H}_{x^1} \otimes \mathcal{H}_{x^2} \otimes \mathcal{H}_{x^3} \Longrightarrow |\psi \rangle = | \psi_{x^1} \rangle \otimes  | \psi_{x^2} \rangle \otimes | \psi_{x^3} \rangle \equiv | \psi_{x^1} \rangle | \psi_{x^2} \rangle | \psi_{x^3} \rangle
\]
Now
\[
\widehat{H} | \psi \rangle = \sum_{i=1}^d \left( \frac{-1}{2m} \left( \frac{ \partial }{ \partial x^i} \right)^2  + \frac{1}{2} m \omega^2 (\widehat{x}^i)^2 \right) |\psi \rangle \equiv \left( \frac{-1}{2m} \left( \frac{ \partial }{ \partial x^i} \right)^2  + \frac{1}{2} m \omega^2 (\widehat{x}^i)^2 \right) |\psi \rangle
\]
The ladder operators make it much more simpler and straightforward than directly solving the differential equation.  Let
\[
\begin{aligned}
  & a_i =  \sqrt{ \frac{m\omega}{2} } \left( \widehat{x}^i + \frac{i}{m\omega} \widehat{p}_i \right) \\
  & a_i^{\dag} = \sqrt{ \frac{m\omega}{2} } \left( \widehat{x}^i - \frac{i}{m\omega} \widehat{p}_i \right)
\end{aligned} \quad \quad \, \begin{aligned}
  & \widehat{x}^i = \sqrt{ \frac{1}{2m\omega} } (a_i + a_i^{\dag}) \\ 
  & \widehat{p}_i = i \sqrt{ \frac{m \omega}{2 } } (a_i^{\dag} - a_i ) \\ 
\end{aligned}
\]
Now let 
\[
N_i = a_i^{\dag} a_i = \frac{m\omega}{2} ((\widehat{x}^i)^2 + \frac{ (\widehat{p}_i)^2}{ (m\omega)^2 } ) - \frac{1}{2}
\]
where the uncertainty principle commutator relation was used: 
\[
\widehat{x}^i \widehat{p}_i - \widehat{p}_i\widehat{x}^i = [\widehat{x}^i , \widehat{p}_i ] =  i 
\]
Thus
\[
\widehat{H} = (N_1 + \frac{1}{2} ) \omega + (N_2 + \frac{1}{2} ) \omega + (N_3 + \frac{1}{2} ) \omega
\]
\[
\widehat{H} | \psi \rangle = ((n_x + \frac{1}{2} )\omega + (n_y + \frac{1}{2}) \omega + (n_z + \frac{1}{2} ) \omega ) | \psi \rangle
\]

Consider 
\[
Z_s = \sum_{n_x,n_y,n_z} \exp{ \left[ - \beta((n_x + \frac{1}{2} )\omega + (n_y + \frac{1}{2}) \omega + (n_z + \frac{1}{2} )\omega -\epsilon_0 ) \right] } = e^{\beta \epsilon_0} \left( \sum_n (e^{-\beta \omega })^n \right)^d = e^{\beta \epsilon_0} \left( \frac{1}{ 1 - e^{-\beta \omega } } \right)^d
\]

The derivation of $p = p(\tau)$ goes, as above (with Example: Model system for gas-solid equilibrium (pp.285 Ch. 10 ``Phase transformations'', Kittel and Kroemer (1980) \cite{CKittelHKroemer1980})); recall
\[
\lambda_g = \lambda_s \, ,  \, \lambda_g = \frac{n}{n_Q} = \frac{p}{\tau n_Q}
\]
and so for $\tau \gg \hbar \omega$ or $1 \gg \beta \omega$, 
\[
p = \left( \frac{M}{2\pi \hbar^2} \right)^{3/2} \tau^{5/2} e^{-\beta \epsilon_0}(1- e^{-\beta \omega} )^d \cong \left( \frac{M}{2\pi } \right)^{3/2} \frac{1}{\hbar^3} \tau^{5/2} e^{-\beta \epsilon_0} \beta^d \omega^d
\]
For $d=3$,
\[
\boxed{ p = \left( \frac{M}{2\pi }\right)^{3/2} \frac{\omega^3}{\tau^{1/2}} e^{-\beta \epsilon_0} }
\]



\item[(b)] 
From above part (a),
\[
\frac{dp}{d\tau}\frac{1}{p} = \frac{-1}{2} \tau^{-1} + (-\epsilon_0 \left( \frac{-1}{\tau^2} \right) )
\]
Now using the approximation that $\Delta v \approx v_g$, and the ideal gas idealization,
\[
L = \frac{dp}{d\tau} \frac{\tau^2}{p} = \epsilon_0 - \frac{\tau}{2}
\]
My explanation: we must excite the atom by its (bounded) ground state energy, $\epsilon_0$.  On the other hand, thermal energy of gas, with increasing linear temperature, helps to remove atom from solid to gas.  
\end{enumerate} 


\problemhead{5} \textbf{Gas-solid equilibrium}

\begin{enumerate}
\item[(a)] 
Now 
\[
F_s = -\tau \ln{ \frac{Z_s^{N_s} }{N_s!} } = -\tau \left( N_s \ln{Z_s} - \ln{N_s!} \right) \simeq -\tau ( N_s \ln{Z_s} - N_s \ln{N_s} + N_s )
\]
where Stirling's approximation was used and $(1/N_s!)$ factor is for indistinguishability (from quantum mechanics).  

Recall
\[
Z_s = \frac{e^{\beta \epsilon_0}}{ (1- e^{-\beta \omega})^d} \text{ so } \ln{Z_s} = \beta \epsilon_0 - d\ln{ (1-e^{-\beta \omega} )}
\]
from Example: Model system for gas-solid equilibrium (pp.285 Ch. 10 ``Phase transformations'', Kittel and Kroemer (1980) \cite{CKittelHKroemer1980}).  

Then
\begin{equation}\label{Eq:Prob1005aF_s}
  F_s = -N_s \epsilon_0 + N_s \tau d\ln{(1-e^{-\beta \omega} )} + \tau N_s \ln{ N_s} - N_s\tau
\end{equation}
Recall this useful relation (it's easy to calculate $\sigma$ from Helmholtz free energy $F$):
\[
\left( \frac{ \partial F}{ \partial \tau} \right)_V = \sigma
\]
\begin{correction}{Helmholtz derivative}
Differentiating $F = U - \tau\sigma$ gives $(\partial F/\partial\tau)_V = -\sigma$.
The minimization in $N_g$ below does not use this line.
The boxed $N_g$ and $p_g$ are unchanged.
See \S\ref{sec:helmholtz-sign}.
\end{correction}
and so
\begin{equation}\label{Eq:Prob1005aentropy}
\begin{aligned}
  \frac{ \partial F_s}{ \partial \tau} & = N_s d\ln{(1- e^{-\omega/\tau} )} + N_s \ln{ N_s} - N_s + N_s \tau d \left( \frac{1}{ 1 - e^{-\beta \omega} } \right)(-e^{-\beta \omega} )(-\omega)\left( \frac{-1}{\tau} \right) = \\
  & = N_s d\ln{ (1- e^{-\beta \omega} )} + N_s \ln{N_s} - N_s - N_s d\beta \omega \frac{e^{-\beta \omega} }{ 1 - e^{-\beta \omega} } = \sigma
\end{aligned}
\end{equation}


Suppose $\tau \gg \omega$ or $1 \gg \beta \omega$.  Then $\sigma = N_s d\ln{ (1-e^{-\beta \omega})} + N_s \ln{N_s} - N_s(1+d)$.  

So take the (given) ``extreme assumption that the entropy of the solid maybe neglected over the temperature range of interest.'' \cite{CKittelHKroemer1980}

Then comparing $\sigma$ from Eq. \ref{Eq:Prob1005aentropy} and $F_s$ from Eq. \ref{Eq:Prob1005aF_s}, then
\[
F_s \simeq -N_s \epsilon_0
\]


\[
\boxed{ F = F_s + F_g = -N_s \epsilon_0 + N\tau \left[ \ln{ \left( \frac{N_g}{Vn_Q} \right) } - 1 \right] }
\]
\item[(b)] Now
\[
F = -N_s \epsilon_0 + N_g \tau [ \ln{ \left( \frac{N_g}{Vn_Q} \right)} -1 ] = -\epsilon_0 (N-N_g) + N_g \tau \left( \ln{ \left( \frac{N_g}{Vn_Q} \right) - 1 } \right)
\]
Now, taking the partial derivative,
\[
\frac{ \partial F}{ \partial N_g} = \epsilon_0 + \tau \left( \ln{ \left( \frac{N_g}{Vn_Q} \right) }-1 \right) + N_g \tau \frac{1}{N_g}
\]
Taking the minimum free energy with respect to $N_g$, i.e. set the partial derivative to $0$,
\[
\frac{ \partial F}{ \partial N_g } = 0 = \epsilon_0 + \tau \ln{\left( \frac{N_g}{Vn_Q} \right) } \text{ or } \boxed{ N_g = n_Q V \exp{ \left( \frac{-\epsilon_0}{\tau} \right) } }
\]
\item[(c)] Idealize the gas with the ideal gas law.  
\[
pV = N_g \tau \Longrightarrow p_g = n_Q \tau \exp{ \left( \frac{-\epsilon_0}{\tau} \right) } = \boxed{ \left( \frac{M}{2\pi \hbar^2 } \right)^{3/2} \tau^{5/2} \exp{ \left( \frac{-\epsilon_0}{\tau} \right) } = p_g }
\]
\end{enumerate}


